\documentclass[../main.tex]{subfiles}

\begin{document}

% ============================================================
\chapter{Empirical Analysis and Model Development}
\label{chap:5-Empiric-analysis-models}
% ============================================================

The experimental evaluation of the impact of controlled data corruption on deep learning models for intrusion detection proceeds through three principal phases: (i)~systematic hyperparameter optimization, (ii)~baseline model creation with feature selection, and (iii)~large-scale noise-injection experiments.  Each phase builds upon the previous, establishing a controlled environment in which the effect of data quality degradation can be rigorously quantified.  This chapter describes the methodology and technical implementation of these three phases in detail.

% ============================================================
\newpage
\section{Phase 1: Hyperparameter Tuning Strategy}
\label{sec:hpo-strategy}
% ============================================================

Modern deep learning architectures exhibit high sensitivity to hyperparameter choices.  While the discovery of optimal hyperparameters remains a computationally expensive endeavor, systematic optimization via machine learning-driven search strategies offers a principled alternative to manual tuning or random search.  This section describes the hyperparameter optimization (HPO) procedure applied to both TabNet and the custom CNN-LSTM architecture.

% ────────────────────────────────────────────────────────────
\subsection{Optimization Framework and Search Strategy}
\label{sec:hpo-framework}
% ────────────────────────────────────────────────────────────

The hyperparameter optimization was conducted using \textbf{Optuna}~\supercite{akiba2019optuna}, an open-source automatic hyperparameter tuning framework that implements sampling and pruning strategies tailored to accelerate the discovery of high-performing configurations.

The search employed the \emph{Tree-structured Parzen Estimator} (TPE) sampler, a Bayesian optimization algorithm that combines kernel-density estimation with a tree-structured data structure to balance exploration and exploitation of the hyperparameter space.  The TPE sampler iteratively constructs probabilistic models of the objective function (in this case, validation-set performance) conditioned on previous trials, thereby directing subsequent evaluations toward regions of the hyperparameter space that are predicted to yield improved performance.

To ensure computational tractability while maintaining search diversity, the HPO campaign was conducted on a \emph{stratified 0.1\% subsample} of the unified CIC-IDS2018 dataset (approximately 16,000 instances stratified by the binary label to preserve class proportions).  This reduced dataset enables rapid iteration through many hyperparameter configurations while remaining representative of the class-label structure of the full dataset.

% ────────────────────────────────────────────────────────────
\subsection{TabNet Hyperparameter Space}
\label{sec:hpo-tabnet}
% ────────────────────────────────────────────────────────────

TabNet's architecture is governed by three principal hyperparameters that govern the sequential attention mechanism and the learning process:

\begin{enumerate}
    \item \textbf{Number of attention steps} ($N_a$): Controls the number of sequential decision steps in the TabNet forward pass.  Each step applies an attention mask to select a sparse subset of input features for processing at that step.  The search space was set to $N_a \in \{4, 6, 8, 10\}$, representing a trade-off between model complexity and feature-selection granularity.

    \item \textbf{Feature dimension size} ($N_d$): Determines the dimension of the linear layers within each decision step.  This hyperparameter indirectly controls the model's capacity to learn complex feature combinations.  The search space was $N_d \in \{32, 64, 128\}$.

    \item \textbf{Sparsity regularization coefficient} ($\lambda_{\text{sparse}}$): A float parameter controlling the weight of the sparsity penalty imposed on the attention masks.  Higher values encourage sparser feature selection (fewer features per step), while lower values allow more features.  The search space was $\lambda_{\text{sparse}} \in [0.000001, 0.01]$ (log-uniform).

    \item \textbf{Learning rate}: The optimization step-size for gradient descent.  The search space was $\text{lr} \in [0.0001, 0.01]$ (log-uniform).

    \item \textbf{Batch size}: The number of instances per gradient-accumulation step.  The search was constrained to batch sizes $\in \{256, 512, 1024\}$, balancing gradient stability with memory consumption.

    \item \textbf{Mask type}: A categorical choice between \texttt{softmax} (probabilistic sparse selection) and \texttt{relu} (hard masking via rectified linear activation).  The TabNet architecture supports both variants; the search evaluated both to determine which is more effective on the NIDS task.
\end{enumerate}

The Optuna trial budget for TabNet was set to 20 trials, allowing thorough exploration of the hyperparameter space within computational limits.

% ────────────────────────────────────────────────────────────
\subsection{CNN-LSTM Hyperparameter Space}
\label{sec:hpo-cnnlstm}
% ────────────────────────────────────────────────────────────

The custom CNN-LSTM architecture combines convolutional layers (designed to extract spatial patterns within time windows) with recurrent layers (designed to model temporal dependencies across windows).  The principal hyperparameters controlling this architecture are:

\begin{enumerate}
    \item \textbf{Number of convolutional filters} ($n_{\text{filters}}$): Controls the width of the convolutional layer (the number of learned convolution kernels).  The search space was $n_{\text{filters}} \in \{16, 32, 64, 128\}$.

    \item \textbf{Convolutional kernel size} ($k_{\text{size}}$): The temporal window covered by each convolution operation.  The search space was $k_{\text{size}} \in \{3, 5, 7\}$.

    \item \textbf{LSTM hidden units (first layer)} ($n_{\text{LSTM1}}$): The dimensionality of the hidden state in the first LSTM layer.  The search space was $n_{\text{LSTM1}} \in \{32, 64, 128, 256\}$.

    \item \textbf{LSTM hidden units (second layer)} ($n_{\text{LSTM2}}$): The dimensionality of the hidden state in the stacked second LSTM layer.  The search space was $n_{\text{LSTM2}} \in \{16, 32, 64, 128\}$.

    \item \textbf{Dropout rate}: The fraction of units randomly deactivated during training to prevent overfitting.  The search space was dropout $\in [0.1, 0.5]$ (continuous uniform).

    \item \textbf{Learning rate}: The optimization step-size, searched over $\text{lr} \in [0.00001, 0.001]$ (log-uniform).

    \item \textbf{Batch size}: Constrained to batch sizes $\in \{256, 512, 1024\}$.
\end{enumerate}

The Optuna trial budget for CNN-LSTM was 20 trials, matching TabNet.

% ────────────────────────────────────────────────────────────
\subsection{Multi-Task Training Configuration}
\label{sec:hpo-multitask}
% ────────────────────────────────────────────────────────────

Both models were instantiated as \emph{multi-task classifiers}: each model simultaneously predicts two labels from the same input:
\begin{itemize}
    \item The \emph{binary label} (\texttt{Label\_generic}): benign vs.\ malign
    \item The \emph{multi-class label} (\texttt{Label}): the specific attack type (or benign)
\end{itemize}

The multi-task setup enables a single model to solve both the anomaly-detection problem (binary classification) and the attack-type identification problem (multi-class classification) within a unified learning framework.  Shared intermediate representations can improve generalization by exploiting commonalities between the two prediction tasks while allowing task-specific output heads to specialize.

Both models employ a \emph{weighted sum} loss combining the two tasks:
\[
\mathcal{L}_{\text{total}} = \alpha \cdot \mathcal{L}_{\text{binary}} + (1-\alpha) \cdot \mathcal{L}_{\text{multi-class}}
]
where $\alpha = 0.5$ each task receives equal weight.  The loss functions are cross-entropy on both branches, with class weights inversely proportional to the training-set frequencies of each class to accommodate the severe class imbalance present in the CIC-IDS2018 dataset.

% ────────────────────────────────────────────────────────────
\subsection{HPO Objective Function and Termination Criteria}
\label{sec:hpo-objective}
% ────────────────────────────────────────────────────────────

Each Optuna trial trains a candidate model on the 0.1\% subsample using a randomly selected hyperparameter configuration drawn from the search spaces described above.  Training proceeds for up to \textbf{20 epochs} (TabNet) or \textbf{50 epochs} (CNN-LSTM), with \emph{early stopping} triggered if the validation loss does not improve for 5 consecutive epochs.

The objective function optimized by Optuna is the \textbf{F1-score} on the multi-class prediction task (attack-type identification), averaged across all attack classes using macro-averaging (unweighted mean of per-class F1 scores).  F1-score was chosen as the primary optimization metric because it balances precision and recall, both critical for an intrusion-detection system: high precision prevents false alarms that exhaust analyst resources, while high recall minimizes missed attacks.

The optimization campaign terminates after 20 trials per model.  Although larger trial budgets might discover incrementally better configurations, 20 trials represent a practical compromise between search quality and computational cost on the 0.1\% subsample.

% ============================================================
\newpage
\section{Phase 2: Feature Importance Analysis and Baseline Model Training}
\label{sec:baseline-creation}
% ============================================================

Following the completion of hyperparameter optimization on the 0.1\% subsample, the best-performing hyperparameter configurations for both TabNet and CNN-LSTM were extracted and used to train baseline models on a larger segment of the dataset.  This phase introduces two principal objectives: (i)~automatic feature importance estimationand ranking, and (ii)~creation of reference \emph{baseline} models against which noise-injection experiments will be compared.

% ────────────────────────────────────────────────────────────
\subsection{Best Hyperparameter Selection and Extraction}
\label{sec:best-hpo-configs}
% ────────────────────────────────────────────────────────────

From the 20 Optuna trials per model, the configuration achieving the highest validation F1-score was selected as the ``best'' hyperparameter set.  These best configurations are then held fixed throughout all subsequent experiments, ensuring that any observed changes in model performance stem from data quality manipulation (via PuckTrick) rather than from hyperparameter re-tuning.

Storing the best hyperparameter configurations enables reproducibility and allows different corruption scenarios to be tested under identical architectural and optimization conditions, isolating the effect of noise from confounding hyperparameter variations.

% ────────────────────────────────────────────────────────────
\subsection{Feature Importance Extraction from Baseline Models}
\label{sec:feature-importance}
% ────────────────────────────────────────────────────────────

One of the distinctive properties of the TabNet architecture is that it provides \emph{native}\ interpretability through its attention mechanism: at each decision step, TabNet learns a sparse mask indicating which input features are selected for processing at that step.  By aggregating these per-step masks across all steps and the model's training trajectory, a \emph{feature-importance score} can be computed for each input feature, representing the model's learned assessment of each feature's contribution to the prediction tasks.

Similarly, the CNN-LSTM model's contribution to feature selection decisions can be inferred from the magnitude of learned weights in its first convolutional layer, although this is less interpretable than TabNet's explicit attention mechanism.

After training the best TabNet and CNN-LSTM configurations on the 0.1\% subsample, feature-importance scores were extracted from both models.  The rationale for this multi-model approach is that combining importance rankings from different architectures may capture complementary aspects of feature informativeness: TabNet's sequential attention reveals which features the model attends to during its step-wise decision process, while the CNN-LSTM's convolutional weights capture which features are deemed salient for temporal pattern matching.

The two importance rankings were then \emph{intersected} to identify features that both models agree are important.  This intersection yielded a \emph{top-K consensus feature set}, where K $\approx 8$--10 features (the exact number depends on ties in the ranking).  This reduced set of ``expert-agreed'' features serves two purposes:
\begin{enumerate}
    \item It reduces dimensionality for subsequent experiments, lowering computational cost.
    \item It reduces the dimensionality of the noise-injection space, focusing experiments on the features most likely to impact model performance.
\end{enumerate}

% ────────────────────────────────────────────────────────────
\subsection{Baseline Model Training on Expanded Dataset}
\label{sec:baseline-training}
% ────────────────────────────────────────────────────────────

With best hyperparameter configurations in hand and the consensus feature set identified, the dataset was scaled up to \textbf{10\% of the full CIC-IDS2018} (approximately 1,617,000 instances).  This 10\% subsample was generated via \emph{stratified sampling} on the binary label to preserve the class proportions of the full dataset.

Critically, the stratified sampling preserves the \emph{temporal sequencing} of the data: rather than randomly shuffling and selecting instances, the sampling strategy selects every $(1/0.1) = 10$-th instance while respecting the temporal order of flow records within the original Parquet file.  This approach ensures that if a malicious campaign occurred as a cluster of consecutive flows in the original dataset, that temporal coherence is preserved in the sampled subset.

Both TabNet and CNN-LSTM models were then re-trained on this 10\% dataset using the previously identified best hyperparameters and the same multi-task training setup.  These models serve as the \emph{baseline} against which all noise-injection experiments will be compared: any performance degradation observed when training on corrupted versions of this dataset directly quantifies the impact of the corruption.

The training protocol is identical to HPO: early stopping is applied with a patience of 5 epochs to avoid overfitting.  The maximum number of epochs was again set to 20 for TabNet and 50 for CNN-LSTM, chosen based on preliminary training curves observed during HPO.

% ────────────────────────────────────────────────────────────
\subsection{Baseline Model Evaluation Metrics}
\label{sec:baseline-metrics}
% ────────────────────────────────────────────────────────────

After training the baseline models, their performance was evaluated on a held-out test set (20\% of the 10\% subsample, approximately 323,400 instances) that remained untouched throughout all experiments.  The evaluation was conducted using a comprehensive set of performance metrics to capture different aspects of model quality:

\begin{itemize}
    \item \textbf{Accuracy}: overall fraction of correctly classified instances (both labels \emph{binary} and \emph{multi-class}).

    \item \textbf{F1-score}: harmonic mean of precision and recall, micro-averaged (aggregated instance-wise) and macro-averaged (unweighted mean of per-class F1 scores).

    \item \textbf{Precision and Recall}: for each label, computed separately to isolate type-I errors (false positives, triggering unnecessary alerts) and type-II errors (false negatives, missed attacks).

    \item \textbf{Matthews Correlation Coefficient (MCC)}: a correlation measure between predicted and actual class labels that is informative even on imbalanced datasets, accounting for all four cells of the confusion matrix.

    \item \textbf{Receiver Operating Characteristic Area Under Curve (ROC-AUC)}: the area under the curve of true-positive rate vs.\ false-positive rate, aggregating performance across all classification thresholds.

    \item \textbf{Confusion matrices}: structured as separate matrices for each label hierarchy (binary and multi-class) to visualize the distribution of correct and incorrect predictions across classes.
\end{itemize}

These baseline performance metrics establish the \emph{clean-data ground truth} against which all corrupted-data performance will be measured.  The difference between clean and corrupted performance, denoted $\Delta \text{performance}$, quantifies the fragility of each model-feature-corruption-level combination.

% ============================================================
\newpage
\section{Phase 3: Systematic Noise-Injection Experiments}
\label{sec:noise-injection-experiments}
% ============================================================

Having established baseline model performance on clean data and identified the most important features for both models, the third phase systematically introduces controlled corruption into selected features and re-trains the models to measure performance degradation.  This controlled experimental setup enables answering the central research question: \emph{how fragile are state-of-the-art deep learning architectures when trained on data corrupted along specific quality dimensions?}

% ────────────────────────────────────────────────────────────
\subsection{Corruption Strategy: Feature and Method Selection}
\label{sec:corruption-strategy}
% ────────────────────────────────────────────────────────────

The noise-injection campaign targets a curated set of features identified in Phase 2 as both \emph{highly important} (high feature-importance scores from both models) and, ideally, \emph{highly fragile} (high skewness and/or kurtosis as characterized in Chapter~\ref{chap:dataset}).  Features were selected from the consensus top-10 list, ensuring that injected corruption is likely to degrade model performance by targeting features that the models rely upon.

The corruption methods employed derive directly from PuckTrick's corruption taxonomy (as described in Section~\ref{sec:pucktrick-limitations} of Chapter~\ref{chap:related}):

\begin{enumerate}
    \item \textbf{Missing values} (via \texttt{missing.py}): randomly replace feature values with \texttt{NaN}.  Missing values force the model to either employ learned imputation strategies or to propagate missing information through the network, both of which can degrade performance.

    \item \textbf{Noisy values} (via \texttt{noisy.py}): replace feature values with random samples drawn uniformly from the observed range $[\min, \max]$ of that feature.  This introduces statistical noise that mimics real-world measurement uncertainty or sensor drift.

    \item \textbf{Outliers} (via \texttt{outliers.py}): inject extreme values beyond the normal distributional range (e.g., values $> \mu + 4\sigma$ for continuous features).  Outliers can disproportionately influence learning if the model is sensitive to extreme values or if the loss function gives large gradients to misclassifications involving outliers.

    \item \textbf{Label corruptions} (via \texttt{labels.py}): flip or reassign the target label for a fraction of instances.  Applied only to the binary and multi-class labels (not to feature columns), label noise directly corrupts the ground truth, forcing the model to learn from partially mislabeled data.

    \item \textbf{Duplicated records} (via \texttt{duplicated.py}): introduce exact duplicate rows into the training set.  Duplicates create artificial oversampling of particular instances, biasing the model towards memorizing those instances.
\end{enumerate}

Applications of methods 1--4 are \emph{feature-specific}: the user specifies a target column and a corruption percentage, and the method corrupts exactly that percentage of instances in that column.  The \emph{duplicated} method, by contrast, operates at the row level: it introduces duplicate rows into the dataset.  In this campaign, duplicates are applied once globally (not per-feature).

% ────────────────────────────────────────────────────────────
\subsection{Experimental Design: Corruption Levels and Combinations}
\label{sec:experimental-design}
% ────────────────────────────────────────────────────────────

For each target feature in the consensus set and each corruption method, a \emph{corruption level grid} was defined specifying the percentage of instances to corrupt.  The grid comprises the following percentages:
\[
P = \{0.05, 0.10, 0.20, 0.35, 0.50, 0.75\}
\]
representing mild (5\%), light (10\%), moderate (20\%), significant (35\%), substantial (50\%), and severe (75\%) corruption levels.

The experimental matrix is thus:
\[
|\text{Experiments}| = |\text{consensus features}| \times |\text{corruption methods}| \times |P|
\]

With approximately 8--10 consensus features, 4 feature-specific methods, and 6 corruption levels, this yields approximately $(8 \text{ to } 10) \times 4 \times 6 = 192 \text{ to } 240$ experiments.  Adding the duplicated method (applied once per seed, independent of features) brings the total to approximately 200--250 experimental runs \emph{per random seed}.

% ────────────────────────────────────────────────────────────
\subsection{Dataset Preparation and PuckTrick Corruption}
\label{sec:pucktrick-workflow}
% ────────────────────────────────────────────────────────────

For each experiment, the following workflow was executed:

\begin{enumerate}
    \item \textbf{Load clean baseline data}: Re-load the 10\% dataset from disk (the same dataset used for baseline training).

    \item \textbf{Apply PuckTrick corruption}: Invoke the appropriate PuckTrick method with the target feature, corruption method, and corruption percentage.  This produces a corrupted variant of the dataset in which the specified fraction of values in the target feature have been substituted according to the corruption semantics.

    \item \textbf{Train on corrupted data}: Using the \emph{same hyperparameters} as the baseline models (determined in Phase 1), train fresh instances of both TabNet and CNN-LSTM on the corrupted dataset.  Training proceeds identically to the baseline phase: multi-task configuration, early stopping, and the same 20 or 50 epoch budgets.

    \item \textbf{Evaluate on clean test set}: Crucially, the test set remains \emph{uncorrupted}.  Model performance is evaluated on the same held-out 20\% test fraction used for baseline evaluation.  This ensures that performance differences reflect the model's ability to generalize from corrupted training data to unseen clean data, a realistic scenario where training data may be imperfect but test data is presumed clean.

    \item \textbf{Record results}: Accuracy, F1-score, precision, recall, MCC, and ROC-AUC are recorded for both models on both label hierarchies (binary and multi-class).  These results are persisted to disk for post-hoc analysis.

\end{enumerate}

% ────────────────────────────────────────────────────────────
\subsection{Statistical Robustness via Repeated Runs with Random Seeds}
\label{sec:robustness-seeds}
% ────────────────────────────────────────────────────────────

To establish statistical confidence in the experimental results and to account for the inherent randomness in deep learning (Gaussian initialization of weights, stochastic gradient descent, dropout), the entire noise-injection experimental loop was repeated \textbf{multiple times with different random seeds}.  Each iteration of the outer loop (one for each random seed) performs all PuckTrick-based experiments using a different seed for:
\begin{itemize}
    \item Dataset splitting (train/validation/test partition generation)
    \item Neural network weight initialization
    \item Stochastic regularization mechanisms (dropout, batch normalization)
    \item Optuna's hyperparameter search (if re-doing HPO per seed)
    \item PuckTrick's corruption sampling (random selection of which instances to corrupt)
\end{itemize}

By repeating the entire campaign 10 times with distinct seeds (e.g., $\{51, 303, 404, 505, 606, 707, 808, 909, 1010, 1111\}$), confidence intervals can be computed around the observed performance degradation metrics.  These confidence intervals account for the uncertainty inherent in the experiment and provide a principled statistical basis for distinguishing genuine effects of corruption from noise.

The seed-based replications also enable the calculation of \emph{robustness statistics}, such as the standard deviation of $\Delta \text{performance}$ across seeds, offering a quantitative measure of how consistently a particular corruption affects model performance.

% ────────────────────────────────────────────────────────────
\subsection{Two Experimental Branches: Reduced vs. Full Feature Set}
\label{sec:two-branches}
% ────────────────────────────────────────────────────────────

Following the strategy outlined in Chapter~\ref{chap:dataset}, experiments are conducted along two complementary branches:

\paragraph{\red{Experiment A: Reduced Feature Set (Proof of Concept)}}
A first set of experiments is conducted using the reduced feature set derived in Section~\ref{sec:reduced-feature-set} of Chapter~\ref{chap:dataset}: the dataset with 44 features (8 zero-variance and 27 high-correlation features removed).  The rationale for this branch is to establish a \emph{proof of concept}: to demonstrate that systematic noise injection, when applied to a curated set of low-redundancy, high-impact features, can either degrade performance (as a null hypothesis) or, in some cases, \emph{improve} performance (if the noise acts as a regularizer).

This branch focuses on identifying noise-injection strategies that consistently improve intrusion-detection performance, a finding that would have immediate practical implications for the development of more robust NIDS systems.

\paragraph{\red{Experiment B: Full Feature Set (Theoretical Depth)}}
In contrast, the second branch conducts experiments on the \emph{full 79-feature dataset} without any feature removal.  This branch addresses the more theoretical aspects of the thesis: understanding the holistic impact of data corruption across the entire feature space.  It enables analysis of how noise propagates through zero-variance and redundant features, how models exploit or ignore such features, and how the presence of feature redundancy affects the overall resilience to corruption.

The two branches produce divergent experimental footprints and lead to distinct conclusions in the final analysis phase (Chapter~\ref{chap:6-conclusions}), as the reduced-feature experiments may reveal promising regularization effects while the full-feature experiments reveal broader patterns of vulnerability and robustness.

% ────────────────────────────────────────────────────────────
\subsection{Computational Infrastructure and Distributed Execution}
\label{sec:computational-infrastructure}
% ────────────────────────────────────────────────────────────

The large-scale noise-injection experiments necessitate significant computational resources.  To manage the workload efficiently, the experimental pipeline was implemented as a scalable Python application that leverages:

\begin{itemize}
    \item \textbf{Apache Spark} (via PySpark) for distributed data loading and PuckTrick corruption operations.  The Spark-extended PuckTrick library (as described in Chapter~\ref{chap:oth}) enables distributed dataset preparation, allowing large datasets to be corrupted in place without collecting all data to a single node.

    \item \textbf{PyTorch and PyTorch Lightning} for distributed multi-GPU training.  While the primary experimental runs have been conducted on single-GPU or CPU environments, the architecture supports seamless scaling to multi-GPU setups via PyTorch Lightning's distributed training utilities.

    \item \textbf{Optuna} within the training loop to manage hyperparameter search results and best-configuration persistence.

    \item \textbf{Checkpoint and artifact persistence}: After each experiment, model checkpoints, training logs, and performance metrics are persisted to local disk (or distributed storage on HPC systems) to enable post-experiment analysis without re-running training.
\end{itemize}

The modular architecture separates \emph{dataset preparation} (executed once per seed, storing results to disk) from \emph{model training} (executed on-demand using pre-prepared datasets).  This decoupling enables:
\begin{enumerate}
    \item Rapid iteration during development without re-executing expensive PuckTrick operations.
    \item Deterministic reproducibility: identical corrupted datasets, when stored, produce identical training behavior given identical random seeds.
    \item Efficient resource allocation: datasets and models can be distributed across compute nodes according to resource availability.
\end{enumerate}

% ========== RESULTS PLACEHOLDER ==========
\subsection{Experiment Execution Summary}
\label{sec:execution-summary}

The full experimental campaign encompasses:
\begin{itemize}
    \item \textbf{Phase 1 (HPO):} 40 hyperparameter optimization trials (20 per model) on a 0.1\% dataset subsample.

    \item \textbf{Phase 2 (Baseline):} 2 baseline model trainings (one per architecture) on a 10\% dataset subsample with fixed best hyperparameters.

    \item \textbf{Phase 3 (Noise Injection):} Approximately 200--250 noise-injection experiments per random seed, conducted across 10 distinct random seeds, for a total of approximately 2,000--2,500 model training runs comprising \red{Experiment A} (reduced features) and \red{Experiment B} (full features).
\end{itemize}

All experiments generate detailed logs, model checkpoints, and performance metrics that form the basis for the analysis presented in the following chapter.

% ============================================================
\newpage
\section{Summary: Experimental Methodology}
\label{sec:methodology-summary}
% ============================================================

Table~\ref{tab:methodology-summary} consolidates the key parameters and design choices of the three-phase empirical investigation:

\begin{table}[htbp]
  \centering
  \caption{Summary of empirical investigation methodology.}
  \label{tab:methodology-summary}
  \begin{tabular}{lll}
    \toprule
    \textbf{Phase} & \textbf{Component} & \textbf{Specification} \\
    \midrule
    \multirow{4}{*}{\textbf{Phase 1: HPO}}
    & Dataset size & 0.1\% ($\approx$16,000 instances) \\
    & Optimization method & Optuna (TPE sampler) \\
    & Trials per model & 20 \\
    & Optimization metric & Macro-averaged F1 (multi-class) \\
    \midrule
    \multirow{5}{*}{\textbf{Phase 2: Baseline}}
    & Dataset size & 10\% ($\approx$1,617,000 instances) \\
    & Models & TabNet + CNN-LSTM \\
    & Training setup & Multi-task (binary + multi-class) \\
    & Epochs & 20 (TabNet), 50 (CNN-LSTM) \\
    & Early stopping patience & 5 epochs \\
    & Test set size & 20\% of 10\% subsample \\
    \midrule
    \multirow{6}{*}{\textbf{Phase 3: Noise}}
    & Target features & 8--10 consensus features \\
    & Corruption methods & 4 feature-specific + 1 row-level \\
    & Corruption percentages & \{0.05, 0.10, 0.20, 0.35, 0.50, 0.75\} \\
    & Experiments per seed & $\approx$200--250 \\
    & Random seeds & 10 distinct seeds \\
    & Total training runs & $\approx$2,000--2,500 \\
    \midrule
    \textbf{Evaluation} & Metrics & Accuracy, F1, Precision, Recall, MCC, ROC-AUC \\
    & Test (clean) & All trained models \\
    & Confidence intervals & Via 10-fold seed-based replication \\
    \bottomrule
  \end{tabular}
\end{table}

This structured, three-phase methodology establishes a rigorous experimental foundation for investigating the interplay between data quality and deep learning performance on a realistic cybersecurity application domain.  The results of these experiments, including performance degradation patterns, model-specific vulnerabilities, and potential noise-driven improvements, are analysed in the following chapter.

\end{document}
